\section{Fazit}
In diesem Experiment wurde die Compton-Streuung untersucht. Als Strahlungsquelle diente eine $^{137}\text{Cs}$-Präparation, deren Gammastrahlen auf Aluminiumtargets unterschiedlicher Dicke gerichtet wurden. Eine zusätzliche $^{152}\text{Eu}$-Quelle wurde zur Kalibrierung der Spektrometerenergie verwendet. Die gestreute bzw. abgeschwächte Strahlung wurde mit einem Szintillationsspektrometer detektiert, das aus einem NaI(Tl)-Kristall und einem Photomultiplier bestand. Die Pulse wurden mithilfe eines Vorverstärkers, eines Hauptverstärkers und eines Vielkanalanalysators (MCA) in Energiehistogramme umgewandelt und anschließend analysiert.

Zunächst wurde die Strahlungsdämpfung mit Aluminiumtargets im Dickenbereich von 1 bis 30 mm untersucht. 
(Abgesehen von einer ungewöhnlichen Beobachtung, bedingt durch ein zu schmales Target, bestätigten die Messdaten den erwarteten exponentiellen Abfall der Transmissionsintensität gemäß dem Lambert-Beer-Gesetz. Der gesamte Streuquerschnitt von Aluminium für Photonen mit einer Energie von ??? keV wurde somit zu $\mu_{tot}$ = ???b bestimmt.)

%Dieser Wert stimmt gut mit dem theoretischen Wert ???????????.

Anschließend wurde das Spektrometer energiekalibriert. Dazu wurde die Cs-Quelle abgeschirmt und eine Eu-Quelle in den Strahlengang eingefügt. Das aufgezeichnete Spektrum zeigte deutlich mehrere Peaks, die den bekannten Übergangsenergien des Europium-Zerfalls entsprachen.
( Eine lineare Regression der Kanalnummern gegen die in der Literatur angegebenen Energien diente als Kalibrierfunktion und bestätigte den linearen Betrieb des Vielkanalanalysators. Durch Vergleich der relativen Intensitäten der Spektrallinien mit bekannten Übergangswahrscheinlichkeiten konnten zudem einige Punkte der relativen Nachweiseffizienz des NaI(Tl)-Detektors bestimmt werden. )
% Aufgrund der begrenzten Anzahl verfügbarer Datenpunkte lieferte die mit einem linearen kontinuierlichen Untergrund angepasste Gauß-Funktion jedoch nur eine eingeschränkte Zuverlässigkeit für die Effizienz über den gesamten Energiebereich.)

In der letzten Messreihe wurde das Compton-Streuspektrum einer $^{137}\text{Cs}$-Quelle mit einem Target fester Dicke von 1 mm und verschiedenen Streuwinkeln aufgenommen. Für jede Messung wurde ein kontinuierliches Hintergrundbild ohne Target erstellt und vom Spektrum subtrahiert. 
(Die Analyse der erhaltenen Spektren ermöglichte eine detaillierte Überprüfung der Compton-Gleichung : Den im Spektrum identifizierten Maxima wurden theoretische Energien zugeordnet und grafisch als Funktion des Streuwinkels dargestellt. Die Anpassung ergab eine einfallende Photonenenergie von: $E_\gamma$ = ???? keV, etwas niedriger/höher/etc als der in der Literatur angegebene Wert von ???? keV.)

In einem weiteren Schritt wurde die Winkelintensität der gestreuten Strahlung mit dem Klein-Nishina-Wirkungsquerschnitt verglichen. Dies erforderte signifikante Korrekturen, insbesondere für die detektorabhängige Effizienz und die vom Targetwinkel abhängige mittlere Dämpfung. 
(Da beide Korrekturen fehlerbehaftet sind, stimmten die erhaltenen Intensitätsprofile (mit und ohne Effizienzkorrektur) nur teilweise mit den theoretischen Werten überein.)

Zusammenfassend lässt sich sagen, dass der gesamte Compton-Streuquerschnitt der $^{137}\text{Cs}$-Quelle quantitativ und die Winkelabhängigkeit der Energie und Intensität der gestreuten Photonen qualitativ bestimmt werden konnten.